% Fragmento de inserción; se incluye desde la raíz de este paquete editorial.
% Las fuentes completas y los cortes materiales constan en ANTECEDENTES_Y_DEPENDENCIAS.md.
% Procedencia: XI REV11 ES, sections/nucleo.tex, tpk_desarrollo_integrado.tex,
% ch_tres_vueltas_ext.tex, ch_elevacion_catalogal.tex, ch_cierre_cardinal.tex,
% ch_descenso_por_fibras.tex; perfil integral public_final/feigenbaum/c37_feigenbaum.tex;
% Grassmannianos REV07 ES, sections/jacobi_spectral.tex;
% TRANSVERSALIDAD_Y_ESCALADO_LOCAL_FEIGENBAUM.md, sección 14.
% Ampliación expositiva de demostraciones existentes; no introduce resultados nuevos.
\providecommand{\ZZ}{\mathbb Z}
\providecommand{\RR}{\mathbb R}
\providecommand{\FF}{\mathbb F}
\providecommand{\APP}{\mathrm{APP}}
\providecommand{\TPK}{\mathrm{TPK}}
\providecommand{\TRIT}{\{-1,0,+1\}}

\section{Operator antecedents of the criticality family}
\label{hmtfg:antecedentes-operatorios}

The criticality family is constructed from a readout of the TPK calendar. The two arithmetic sheets, tritic orientation and transport with memory first determine the states and their prolongations. The ordered reading of the continuum then allows analytic functions to be evaluated, parameters to be compared and roots to be selected. This development preserves the genealogical state and its two readouts: the Archimedean chart provides the evaluation space, while the incidence readout preserves boundary relations. The joint discrete structure of the continuum retains its five consubstantial constructions. The present exposition uses its ordered reader and preserves the common antecedent from which it originates.

\subsection{Arithmetic cells, orientation and transport}
\label{hmtfg:app-trit-transporte}

% XI REV11, sections/nucleo.tex: definiciones de APP y TRIT;
% tpk_desarrollo_integrado.tex:22--165, transporte levantado y calendario.
Let \(D_9=\{1,\ldots,9\}\) and \(X_9=(\ZZ/9\ZZ)^2\). For
\(n\in\ZZ\), division with a positive representative is
\[
 \rho_9(n)=1+((n-1)\bmod9),\qquad
 q_9(n)=\frac{n-\rho_9(n)}9,\qquad
 n=\rho_9(n)+9q_9(n).
\]
The APP cell associated with \((p,q)\in D_9^2\) is
\[
 a_{p,q}=(p,q;R^\Sigma,q^\Sigma,R^\Pi,q^\Pi),
 \quad
 p+q=R^\Sigma+9q^\Sigma,\qquad
 pq=R^\Pi+9q^\Pi.
\]
The two pairs \((R^\Sigma,q^\Sigma)\) and \((R^\Pi,q^\Pi)\) retain the additive and multiplicative evaluations, respectively. The division formula reconstructs both evaluations before any subsequent reduction. Selection of an active sheet is recorded in a state coordinate and preserves the complete cell.

The balanced tritic representative \(r_3(n)\in\{-1,0,+1\}\) satisfies
\(n\equiv r_3(n)\pmod3\). Its carry is
\[
 \chi(a,b)=\frac{a+b-r_3(a+b)}3,\qquad
 a,b\in\{-1,0,+1\}.
\]
The equality
\[
 \chi(a,b)+\chi(r_3(a+b),c)
 =\chi(b,c)+\chi(a,r_3(b+c))
\]
is obtained by writing both sides as
\((a+b+c-r_3(a+b+c))/3\). This identity preserves the carry when regrouping a tritic sum. The operative selector is
\[
 M_{\rm ph}(r)=
 \begin{cases}
 +1,&r\in\{1,4,7\},\\
 -1,&r\in\{2,5,8\},\\
 0,&r\in\{3,6,9\}.
 \end{cases}
\]
The value \(+1\) activates the additive sheet, the value \(-1\) the multiplicative sheet, and the value \(0\) preserves the position of both cursors during the neutral event. All three events belong to the complete chronology and are recorded in the route.

Let \(\mathcal D=\{N,E,S,O\}\), with vectors
\(\nu_N=(-1,0)\), \(\nu_E=(0,1)\), \(\nu_S=(1,0)\) and
\(\nu_O=(0,-1)\), and \(\Theta_{108}=\ZZ/108\ZZ\). The observable configuration is
\[
 q=(x_+,d_+,x_\times,d_\times,\theta)
 \in\mathcal Q_{\rm obs}=(X_9\times\mathcal D)^2\times\Theta_{108}.
\]
For the representative \(\bar\theta\in\{0,\ldots,107\}\), define
\[
 r(\theta)=1+(\bar\theta\bmod9),\qquad
 \beta(\theta)=\left[\left\lfloor\frac{\bar\theta}{9}\right\rfloor\right]_{12},
\]
\[
 a_+(\theta)=\mathbf1_{\{1,4,7\}}(r(\theta)),\qquad
 a_\times(\theta)=\mathbf1_{\{2,5,8\}}(r(\theta)).
\]
The position transition is
\[
 x_+'=x_++a_+(\theta)\nu_{d_+},\qquad
 x_\times'=x_\times+a_\times(\theta)\nu_{d_\times}.
\]
The involution \(\jmath\) exchanges \(N,S\) and \(E,O\). With
\[
 h(\theta)=\mathbf1_{\{0,27,54,81\}}\bigl((\bar\theta+1)\bmod108\bigr),
\]
the complete observable operator is fixed by
\[
 F_{\rm obs}(q)=
 (x_+',\jmath^{h(\theta)}d_+,x_\times',
   \jmath^{h(\theta)}d_\times,\theta+1).
\]

\begin{proposition}[Conservation of lifted transport and observable return]
\label{hmtfg:transporte-calendario}
Transport preserves the integer displacement through the record of boundary crossings. The operator \(F_{\rm obs}\) is invertible and satisfies
\(F_{\rm obs}^{108}=I\).
\end{proposition}
\begin{proof}
Let \(s:X_9\to D_9^2\) be the section of positive representatives. For
\(x'=x+a\nu_d\), \(a\in\{0,1\}\), define
\[
 b_d(x,a)=\frac{s(x)+a\nu_d-s(x')}{9}\in\ZZ^2.
\]
If \(z'=z+b_d(x,a)\), then
\[
 s(x')+9z'=s(x)+9z+a\nu_d.
\]
Summing this equality along a route gives
\[
 s(x_n)+9z_n=s(x_0)+9z_0+\sum_{t=0}^{n-1}a_t\nu_{d_t}.
\]
Partitioning the sum preserves the identity under concatenation. To invert \(F_{\rm obs}\), first subtract one unit from the phase; that phase determines the reversal of directions and the displaced cursor. The previous directions are recovered and the corresponding vector is subtracted. Each block of \(27\) steps contains nine activations of each cursor. The next block reverses directions and cancels their integer displacements. In \(54\) steps, positions and directions are restored; in \(108\), the calendar coordinate is also restored.
\end{proof}

In addition to this configuration, the complete state preserves both arithmetic evaluations, carries, orientation, trajectory, prefix, compatible cylinder, its boundary, memory degree and survival record. Its update factors as
\[
 U_t=\mathsf{Mem}_t\circ\mathsf{Tra}_t\circ\mathsf{Sel}_t.
\]
The holonomy \(\Gamma_9\) composes nine complete transitions. Phase return and memory advance are different coordinates of this composition. The following construction fixes the transitions and their update field by field.

% Recuperación íntegra del propietario de 855 líneas, con etiquetas prefijadas.
\input{\HMTFGRoot/latex/00a_transiciones_nonadicas.tex}

\subsection{Cylinder refinement and ordered readout}
\label{hmtfg:refinamiento-ordenado}

% XI REV11, ch_elevacion_catalogal.tex: bloques de 54 aristas, beta_blk,
% partición efectiva de 729 subcilindros; ch_cierre_cardinal.tex:375--497.
For \(\mathbf b=(b_1,\ldots,b_6)\in\FF_3^6\), let
\(\upsilon:\FF_3\to\{-1,0,+1\}\) be the balanced representative. The three prolongations constructed edge by edge define
\[
 \Gamma^{[54]}_{\mathbf b,k}
 =\gamma_{\upsilon(b_6),k+45}\circ\cdots\circ
   \gamma_{\upsilon(b_1),k}.
\]
Starting from \(Y_0=\{x_\star\}\), put
\[
 Y_{N+1}
 =\{\Gamma^{[54]}_{\mathbf b,1+54N}(x):
          x\in Y_N,\ \mathbf b\in\FF_3^6\}.
\]
Truncation \(\rho^{\rm blk}_{N+1,N}\) removes the last fifty-four edges, recovering the previous state through the transition record.

\begin{proposition}[Effective coding of prolongation blocks]
\label{hmtfg:bloques-729}
Reading the six tritic pairs of each block defines bijections
\[
 \beta_N:Y_N\xrightarrow{\ \cong\ }(\FF_3^6)^N
\]
that commute with truncation. Each state of \(Y_N\) has exactly \(729\) block prolongations in \(Y_{N+1}\).
\end{proposition}
\begin{proof}
On each turn, the record retains the tritic pair and its carry. The three pairs produce the three distinct outputs \(\chi(\varepsilon,\varepsilon)=\varepsilon\), and the record allows recovery of the branch used, including the zero-carry branch. The six concatenations therefore produce all words of \(\FF_3^6\). Two different words are distinguished at the first recorded pair in which they differ. Recovery of the prefix proves compatibility with truncation. By induction on \(N\), every word of length \(N\) has a unique antecedent state; the next word can be any of the \(3^6=729\) possibilities.
\end{proof}

Let \(\Omega_{\Gamma}^{\rm cal}=\varprojlim_NY_N\). The preceding bijections induce
\[
 \Omega_{\Gamma}^{\rm cal}\cong(\FF_3^6)^{\mathbb N}.
\]
The correspondence and its inverse are continuous: each finite coordinate depends on a finite prefix. The cylinders are open and closed. Each cylinder of depth \(N\) is partitioned into its \(729\) subcylinders of depth \(N+1\); each contains an infinite prolongation, obtained, for example, by continuing with the zero branch in all subsequent blocks.

% Antecedentes efectivos del continuo conjunto y del universo interno.
% Se imprimen en ambos destinos; la etiqueta de universo no reemplaza su regla.
\input{\HMTFGRoot/latex/00b_continuo_universo.tex}

\subsection{Ordered readout of block cylinders}
\label{hmtfg:publicacion-cilindros-bloque}

For the representative \([b]_3\in\{0,1,2\}\), fix
\[
 \nu(\mathbf b)=\sum_{i=1}^6[b_i]_3\,3^{6-i}\in\{0,\ldots,728\}.
\]
In the HMT ordered completion, a prefix
\(s=(m;\mathbf b_0,\ldots,\mathbf b_{N-1})\), \(m\in\ZZ\), determines
\[
 a_N(s)=m+\sum_{j=0}^{N-1}\nu(\mathbf b_j)\,729^{-j-1},\qquad
 I_N(s)=[a_N(s),a_N(s)+729^{-N})_{\rm HMT}.
\]
The carry preserves the joining of the two representatives of a boundary. If \(C_N(s)=\overline{I_N(s)}\), the preceding relations give
\[
 I_N(s)=\bigsqcup_{\mathbf b\in\FF_3^6}I_{N+1}(s;\mathbf b),
 \qquad \operatorname{diam}C_N(s)=729^{-N}\longrightarrow0.
\]

\begin{proposition}[Coverage of the Archimedean readout]
\label{hmtfg:cobertura-arquimediana}
The map
\[
 P_{\rm ar}:\ZZ\times\Omega_\Gamma^{\rm cal}
       \longrightarrow\mathbb R_{\rm HMT},\qquad
 \{P_{\rm ar}(m,\xi)\}
    =\bigcap_N C_N(m;\beta_N(\rho_{\infty,N}\xi))
\]
is well defined, is surjective and determines
\[
 (\ZZ\times\Omega_\Gamma^{\rm cal})/\!\sim_{P_{\rm ar}}
       \ \cong\ \mathbb R_{\rm HMT}.
\]
\end{proposition}
\begin{proof}
The left endpoints form a Cauchy sequence in the HMT completion, since the intervals are nested and their diameters tend to zero. Their limit belongs to all the closed intervals and is unique. For any element of the completion, first choose the integer interval containing it and then the unique half-open subinterval of each partition containing it. The preceding proposition realizes each of those choices through an effective TPK prolongation. The compatible prefixes define \(\xi\), whose readout is the chosen element. At boundaries, the carry identifies equivalent representations after their construction. By definition and surjectivity, the quotient by equality of readout is in bijection with the completion.
\end{proof}

The construction is also carried out in the internal genealogical universe \(\mathfrak G=\mathfrak G_{\rm HMT}(h,m_\infty)\). Here \(h\) encodes the operator rules and \(m_\infty\) the compatible history; the universe is obtained through initial transitive closure, definability at each successor and union at limit stages:
\[
 G_0=\operatorname{TC}\{\omega,u_{h,m},h,m_\infty\},\qquad
 G_{\beta+1}=\operatorname{Def}_{\Sigma_{\rm HMT}}(G_\beta),\qquad
 G_\lambda=\bigcup_{\beta<\lambda}G_\beta .
\]
The relations of \(\Sigma_{\rm HMT}\) are those of the operations already constructed and their codes. The internal domain is therefore retained in all coverage claims. Denote its completion by \(\mathbb R_{\rm HMT}^{\mathfrak G}\) and the ordered realization within the same universe by \(\mathbb R^{\mathfrak G}\).

\begin{proposition}[Internal ordered chart]
\label{hmtfg:carta-interna}
There exists a unique isomorphism of complete Archimedean ordered fields
\[
 \iota:\mathbb R_{\rm HMT}^{\mathfrak G}
             \xrightarrow{\ \cong\ }\mathbb R^{\mathfrak G}
\]
that fixes the unit. It preserves rational endpoints, positive square roots and limits of convergent sequences.
\end{proposition}
\begin{proof}
The unit fixes the integers and rationals. Each element of the HMT completion is associated with the cut of rationals less than it; completeness of the ordered realization provides the unique element with that cut. Rational density proves injectivity and strict order preservation. The same procedure in the reverse direction proves surjectivity. Rational approximations from above and below preserve addition and multiplication and thereby the field structure. The positive element whose square is \(a>0\) is transported to the positive element whose square is \(\iota(a)\). If a sequence converges, for every positive rational tolerance its final terms lie in the rational neighborhood of the limit; order preservation transports that condition. All these operations are performed within \(\mathfrak G\). Uniqueness follows from preservation of rational cuts.
\end{proof}

\subsection{Descent along fibers and preservation of the selector}
\label{hmtfg:descenso-selector}

% XI REV11, ch_descenso_por_fibras.tex:10--39; prueba completa pertinente.
\begin{proposition}[Descent criterion for an operation]
\label{hmtfg:criterio-descenso}
Let \(q:\Omega\twoheadrightarrow J\) and
\(\star:D\subseteq\Omega^2\to\Omega\), with \(D\) saturated by the fibers of \(q\times q\). There exists a unique operation \(\bar\star\) on \((q\times q)(D)\) such that
\[
 q(x\star y)=q(x)\,\bar\star\,q(y)
\]
if and only if
\[
 q(x)=q(x'),\quad q(y)=q(y')
 \ \Longrightarrow\ q(x\star y)=q(x'\star y')
\]
for pairs in the domain.
\end{proposition}
\begin{proof}
If \(\bar\star\) exists, both results are the evaluation of the same pair of arguments. Conversely, choose antecedents of the visible arguments and define the result as the readout of their operation. Saturation fixes the domain; the condition on fibers makes the result independent of the chosen antecedents. Surjectivity proves uniqueness.
\end{proof}

In an application through \(P_{\rm ar}\), this criterion identifies the precise condition for transporting an operation from complete states to their readout. In the chart \(\iota\), the field structure and limits already satisfy that condition. Next, the analytic profile is constructed from the calendar and this compatibility is used to transport its iteration and roots.

\subsection{Dodecaphase readout and quadratic normalization}
\label{hmtfg:perfil-dodecafasico}

% Propietario integral c37_feigenbaum.tex: comienzo hasta la familia f_lambda.
The oriented criticality section on the twelve calendar windows is
\[
 \mathcal P_{\rm crit}(m)=\varphi_m
   =(30m-10)\frac{\pi_{\rm HMT}}{180}
   =\pi_{\rm HMT}\frac{3m-1}{18},\qquad m=1,\ldots,12.
\]
This section fixes the origin and orientation of its angular representatives. The uniform reader of the twelve windows is
\[
 \eta_0(F)=\frac1{12}\sum_{m=1}^{12}F(\varphi_m).
\]
The function \(\eta_0\), which assigns weights to these windows, is distinguished from the ordered chart \(\iota\). Uniform normalization and the angular section are part of the operator data of this profile.

Setting \(a_m=(3m-1)/18\), one obtains
\[
 D_0=\frac1{12}\sum_{m=1}^{12}a_m^2=\frac{899}{648},
 \qquad
 s_0=\frac{36}{\pi_{\rm HMT}\sqrt{899}},\qquad
 k_m=s_0\varphi_m=\frac{2(3m-1)}{\sqrt{899}}.
\]
Indeed, \(\sum_{m=1}^{12}(3m-1)^2=5394=6\cdot899\), so that \(s_0^2\pi_{\rm HMT}^2D_0=2\). The profile is
\begin{equation}
 u_0(x)=\frac1{12}\sum_{m=1}^{12}\bigl[1-\cos(k_mx)\bigr],
 \qquad f_\lambda(x)=1-\lambda u_0(x).
 \label{hmtfg:perfil-familia}
\end{equation}
The constant \(\pi_{\rm HMT}\) enters as an angular readout of the HMT kernel and cancels algebraically in the normalization of the profile. The scale \(\lambda\) is the family parameter, whose critical value will be selected through the dynamics.

\begin{proposition}[Normalization and closed form of the profile]
\label{hmtfg:forma-perfil}
The profile \(u_0\) is even and analytic, satisfies
\(u_0(0)=u_0'(0)=0\) and \(u_0''(0)=2\), and, for
\(y=x/\sqrt{899}\), admits the expression
\[
 u_0(x)=1-\frac{\cos(37y)\sin(36y)}{12\sin(3y)}
       =1-\frac{\sin(73y)-\sin(y)}{24\sin(3y)}.
\]
The apparent singularities are interpreted by continuity. Its initial expansion is
\[
 u_0(x)=x^2-\frac{715769}{2424603}x^4+O(x^6).
\]
\end{proposition}
\begin{proof}
The finite cosine sum preserves evenness and analyticity. Its second derivative at zero is \(12^{-1}\sum_m k_m^2=2\). The sum of the trigonometric progression \(\sum_{m=1}^{12}\cos((6m-2)y)\) is \(\cos(37y)\sin(36y)/\sin(3y)\); the product--sum identity provides the second expression. The cosine series and \(\sum_m(3m-1)^4=4294614\) give the stated fourth-degree coefficient. The expression as a finite sum fixes the extension at the zeros of the closed form’s denominator.
\end{proof}

\subsection{Transport of iterates, roots and minima}
\label{hmtfg:transporte-raices}

% Reunión existente: TRANSVERSALIDAD_Y_ESCALADO_LOCAL_FEIGENBAUM.md, sección 14.1.
Within \(\mathbb R_{\rm HMT}^{\mathfrak G}\), the convergent series
\[
 \cos_{\rm HMT}(t)=\sum_{j=0}^{\infty}
                   \frac{(-1)^j t^{2j}}{(2j)!}
\]
defines the corresponding analytic realization. Let
\[
 u_{\rm HMT}(x)=\frac1{12}\sum_{m=1}^{12}
  \left[1-\cos_{\rm HMT}
   \left(\frac{2(3m-1)x}{\sqrt{899}_{\rm HMT}}\right)\right],
 \qquad F_a(x)=1-a\,u_{\rm HMT}(x),
\]
\[
 S_n^{\rm HMT}(a)=F_a^{\,2^n}(0),\qquad
 S_n(\lambda)=f_\lambda^{\,2^n}(0).
\]

\begin{proposition}[Ordered conjugacy of the family and the root selector]
\label{hmtfg:conjugacion-selector}
For the internal arguments of the preceding expressions,
\[
 \iota\circ F_a=f_{\iota(a)}\circ\iota,\qquad
 \iota(S_n^{\rm HMT}(a))=S_n(\iota(a)).
\]
The chart transports exactly the conditions of return, exact period and parameter order. If a set of admissible roots has a minimum, the image of that minimum is the minimum of the image set.
\end{proposition}
\begin{proof}
The chart fixes the rationals, transports the positive square root and preserves limits of the cosine partial sums. Hence, \(\iota(u_{\rm HMT}(x))=u_0(\iota(x))\). Compatibility with addition and multiplication proves the first equality; induction on the number of iterations proves the second. A root is characterized by equality to zero, which is preserved in both directions. The exact period \(2^n\) adds the inequalities \(F_a^{2^j}(0)\ne0\), \(0\le j<n\), also preserved by the injectivity of \(\iota\). If the parameter is required to exceed a preceding center, that inequality is transported by order preservation. Finally, \(a\le b\) for every \(b\) in the set is equivalent to \(\iota(a)\le\iota(b)\) for every element of its image. Surjectivity onto the sets thus defined completes the assertion about the minimum.
\end{proof}

The domain of this transport is the indicated internal continuum. The existence of each minimum and its identification with a specific dynamical branch are results of the criticality construction that follows; the chart preserves those properties once they have been established.

\subsection{Finite spectral representation of the uniform reader}
\label{hmtfg:jacobi-finito}

% Grassmannianos REV07, sections/jacobi_spectral.tex:35--136,
% mecanismo de Hankel, Gram--Schmidt, recurrencia y representación espectral.
% Especialización atómica ya reunida en la sección 14.2 del desarrollo Feigenbaum.
The squared frequencies of the profile define the positive measure
\[
 s_m=k_m^2=\frac{4(3m-1)^2}{899},\qquad
 \nu_0=\frac1{12}\sum_{m=1}^{12}\delta_{s_m},\qquad
 M_r=\int s^r\,d\nu_0(s).
\]
Its twelve nodes are distinct and positive. For \(1\le r\le12\), the Hankel matrix \(H_r=(M_{i+j})_{0\le i,j<r}\) is positive definite. Indeed, for \(p\) of degree less than \(r\),
\[
 \mathbf p^{\mathsf T}H_r\mathbf p
   =\frac1{12}\sum_{m=1}^{12}p(s_m)^2>0
\]
unless \(p=0\), because a polynomial of degree at most eleven that vanishes at twelve distinct nodes is zero. At degree twelve appears the polynomial \(\prod_m(s-s_m)\), with zero norm in this measure. The corresponding spectral representation has dimension twelve.

Let \(p_0,\ldots,p_{11}\) be the monic orthogonal polynomials,
\(h_j=\int p_j^2\,d\nu_0>0\) and
\(\psi_j=p_j/\sqrt{h_j}\). Since \(\nu_0\) is a probability measure,
\(\psi_0=1\). Multiplication by \(s\) in \(L^2(\nu_0)\) is self-adjoint and has a tridiagonal matrix \(J_{12}\) in this basis.

\begin{proposition}[Exact recurrence and spectral intertwining]
\label{hmtfg:jacobi-entrelazamiento}
With
\[
 a_j=\frac{\int s p_j(s)^2\,d\nu_0(s)}{h_j},\qquad
 \beta_j=\frac{h_j}{h_{j-1}}>0\quad(1\le j\le11),
\]
the matrix \(J_{12}\) has diagonal \(a_j\) and subdiagonal
\(\sqrt{\beta_j}\). If
\[
 Q_{mj}=\frac{\psi_j(s_m)}{\sqrt{12}}
 \quad(1\le m\le12,\ 0\le j\le11),\qquad
 D=\operatorname{diag}(s_1,\ldots,s_{12}),
\]
then
\[
 Q^{\mathsf T}Q=I,\qquad DQ=QJ_{12},\qquad
 J_{12}=Q^{\mathsf T}DQ,
\]
and the profile satisfies
\begin{equation}
 u_0(x)=e_0^{\mathsf T}
          \bigl[I-\cos(x\sqrt{J_{12}})\bigr]e_0.
 \label{hmtfg:perfil-espectral}
\end{equation}
\end{proposition}
\begin{proof}
Positivity of the Hankel systems uniquely provides the monic orthogonal polynomials. On expanding \(s p_j\) in the orthogonal basis, the coefficients of \(p_i\), \(i\le j-2\), vanish because
\(\langle sp_j,p_i\rangle=\langle p_j,sp_i\rangle=0\).
The coefficient of \(p_{j+1}\) is one by monicity, that of \(p_j\) is \(a_j\) and that of \(p_{j-1}\) is \(h_j/h_{j-1}\). At the last degree, the same equality is understood in \(L^2(\nu_0)\), where the monic polynomial of degree twelve vanishing at the nodes represents the zero vector. Normalization gives the stated positive subdiagonal.

Orthonormality is equivalent to \(Q^{\mathsf T}Q=I\). The recurrence equality at each node gives \(DQ=QJ_{12}\). Since \(Q\) is square, we also have \(QQ^{\mathsf T}=I\), and obtain the diagonalization of \(J_{12}\). Its positive eigenvalues allow its positive square root to be defined. For every function defined at those twelve nodes,
\[
 e_0^{\mathsf T}\Phi(J_{12})e_0
   =\frac1{12}\sum_{m=1}^{12}\Phi(s_m),
\]
since \(Qe_0=(1,\ldots,1)^{\mathsf T}/\sqrt{12}\).
The choice \(\Phi(s)=1-\cos(x\sqrt{s})\) proves
\eqref{hmtfg:perfil-espectral}.
\end{proof}

The representation preserves exactly the uniform functional, its moments and its profile. Its first coefficients are
\[
 a_0=M_1=2,\qquad
 \beta_1=M_2-M_1^2=\frac{2493348}{808201}.
\]
The complete genealogy remains in the TPK state from which the reader originates; \(J_{12}\) represents its specific spectral readout. The Gram--Schmidt construction and the Jacobi recurrence are applied to this atomic measure, terminating in dimension twelve.

The operator chain used in the following sections is thus fixed: APP cells and TRIT orientation feed the TPK transition; its prolongations produce cylinders and their ordered readout; the twelve oriented windows and the uniform functional produce \(u_0\); the family \(f_\lambda\) determines returns and roots; the internal chart preserves the order of their selectors. The spectral realization describes the same profile through a finite self-adjoint matrix and a cyclic vector.
